\section{What does entry into the feast give?}
\zlabel{triptych:concise:commentary:start}
The banquet begins with the king's provision, not the guests' achievement. The feast is ready when the servants go out. Refusal does not expose a shortage in the host's resources; it reveals the invited person's unwillingness to enter the offered relation. Matthew's farm and business are ordinary activities made into reasons for neglect. Gregory stresses absorption in earthly work and gain; Chrysostom stresses the negligent heart behind the excuses. Neither requires us to condemn work or possessions simply for existing. The question is what has become more desirable than the call.\footnote{Gregory the Great, \work{Homiliae in Evangelia} 38.5; John Chrysostom, \work{Homilies on Matthew} LXIX.1. Chrysostom also invokes Luke's banquet details; Matthew here names farm and business.}

Paul's experience makes that question more demanding. Abundance itself requires training. Chrysostom explains that plenty can encourage pride or idleness just as want can tempt a person toward wrongdoing. Contentment is learned fidelity in changing circumstances, supported by Christ. It is therefore neither a temperament that happens to want little nor an immunity to bodily needs. The guest with a thriving business and the apostle with an empty stomach face different circumstances, but each must receive the good God gives without making circumstances master of the will.\footnote{Chrysostom, \work{Homilies on Philippians} XV, ad 4:11--13.}

Gregory approaches the new guests from another direction. Bad and good enter the room together. Since a guest can subsequently be expelled in the long Gospel, he identifies Matthew's feast with the present Church. Completed heavenly happiness would not be followed by that loss. This inference places judgment within the hearer's own community. The warning cannot be confined to the original refusers while the present guests claim unconditional security. Entry is real, but entry and final fulfillment are different moments.\footnote{Gregory, \work{Homiliae} 38.1, 7--8.}

That distinction is compatible with Thomas's account of Isaiah's feast. Thomas expressly connects Isaiah with Matthew's royal wedding, reads the mountain through Christ's Passion and judgment, and distinguishes the banquets of the Church, the soul and heaven. The Church receives nourishment from the saving Passion; the soul grows in love and contemplation; heaven brings fullness. Gregory asks how the present assembly can include someone not ready for its final joy. Thomas asks how one promised abundance can already be participated in without being exhausted. Their answers keep present gift and final communion together without making them identical.\footnote{Thomas Aquinas, \work{Expositio super Isaiam ad litteram}, cap. 25, paragraph 84985, exposition of \textit{Et faciet Dominus} and final \textit{Nota}.}

The distinction matters for both welcome and discipline. A mixed gathering does not prove that the invitation was a mistake; the servants gathered as commanded. Nor does another guest's fault give the righteous permission to abandon charity. Gregory's grain among chaff and roses among thorns make life together a test of the one who notices what is wrong. Love need not approve wrongdoing. It must nevertheless seek the good of the person who does it. The feast's generosity reaches those needing conversion, while its holiness makes that conversion consequential.

The short Gospel ends with this mixed company. It retains the rejected invitation, the burning of the city and the new gathering; it does not proclaim the garment or the final called/chosen saying. Its immediate question is whether the hearer answers the call and receives those gathered with him. The sharper distinction between entry and a guest's inspection comes from the long form. A garment-centered explanation consequently draws on a scene that this shorter proclamation omits.

The city's burning also receives different explanations. Chrysostom connects it with Rome's destruction of Jerusalem; Gregory interprets the city's punishment through body and soul under judgment, with angelic ministers as the king's armies. Agreement about conversion does not erase that difference. Jesus' immediate conflict concerns the leaders addressed in the Temple dispute; the long form then judges someone inside the new gathering. Ancient collective polemic does not establish inherited ethnic guilt or authorize violence against a people. The Christian hearer remains among those addressed by the warning.\footnote{Chrysostom, \work{Matthew} LXIX.1; Gregory, \work{Homiliae} 38.5; Mt 21:23, 45; 22:11--14.}

\section{Why is love required when mercy is free?}
Gregory identifies the garment by asking what an admitted guest can still lack. Baptism or faith, in this interpretation, already accounts for entry; neither by itself explains the missing garment. Charity does. Love brought the Son into union with humanity, so love belongs to those celebrating his wedding. The garment expresses the feast's reality rather than an arbitrary test unrelated to it. Refusing love from inside contradicts the communion for which one has been invited.\footnote{Gregory, \work{Homiliae} 38.9--10.}

Chrysostom's account is compatible but differently focused. Calling and cleansing come through grace; the ensuing life demands the recipient's care. His garment is life and practice. Gregory identifies the inner form of the response, Chrysostom its lived expression. Neither makes good conduct a price paid before the original invitation. Both insist that the gift changes what the guest is called to become. The silent guest has no defense before the one who knows him; neither exposition depends on a story about garments distributed at the door, a custom not established by the inspected evidence.\footnote{Chrysostom, \work{Matthew} LXIX.2; Gregory, \work{Homiliae} 38.9, 12.}

Augustine's Entrance psalm supplies the indispensable beginning of that change. Before strict accounting, no conscience can safely claim innocence. Forgiveness and Christ's redemption permit the sinner to approach. Yet mercy in Augustine's exposition also warns concerning the life ahead and accompanies the forgiven person along the way. The garment's demand addresses someone already dependent on pardon. It gives no ground for self-congratulation over those still struggling to enter.\footnote{Augustine, \work{Enarrationes in Psalmos} 129.2--3, 7--8; NPNF1 VIII, Psalm CXXX.}

The Collect's subject, grace before and after good action, joins beginning to perseverance. Thomas makes a related distinction when expounding the Shepherd psalm: mercy inspires the heart and helps complete the work. Augustine's Johannine exposition likewise says that the believer purifies himself with God's assistance. The person is neither self-redeeming nor inert. A request for help makes sense precisely because a real act of love remains to be done.\footnote{Thomas, \work{Super Psalmo} 22, n. 2; Augustine, \work{Homilies on the First Epistle of John} IV.7, on 1 Jn 3:3, the context beyond the Communion excerpt.}

Gregory's twofold weave gives this act its measure: love of God and love of neighbor belong to one garment. Contemplation cannot excuse indifference to a neighbor, and service cannot replace the God toward whom life tends. The test extends beyond mutual affection. Loving a friendly person may be reciprocity; charity reaches the enemy for God's sake. Envy, hatred or an intention secretly to harm another can survive beneath a respectable place at the table. Gregory's account searches the established guest's heart, not merely his outward membership.\footnote{Gregory, \work{Homiliae} 38.10--11.}

The material gift in Philippians prevents that search from becoming exclusively inward. Paul praises the donors for sharing his tribulation. Their care has taken bodily form, just as Gregory's account of judgment names hands that gave nothing and feet that declined to visit the sick. A person can bind those powers by refusing to use them in love. The long Gospel's judgment makes that chosen bondage manifest. Charity's inner reality thus appears in going, giving and serving, rather than in a favorable feeling about one's own intentions.\footnote{Gregory, \work{Homiliae} 38.13; Phil 4:14; Chrysostom, \work{Philippians} XV.}

This requirement does not condemn every imperfect act of love as absence of the garment. Gregory explicitly distinguishes possessing it imperfectly from lacking it, and directs the imperfect toward hope in the King's mercy. His warning against presumption concerns perseverance, not a demand that the struggling person despair. Augustine's account of desire gives growth its positive shape: the heart becomes more capable of what it hopes to receive. Serious conversion and confidence in mercy can therefore belong to the same person.\footnote{Gregory, \work{Homiliae} 38.12, 14; Augustine, \work{First Epistle of John} IV.5--7.}

\section{What does abundance mean for someone still hungry?}
The psalm Communion option promises good to those who seek the Lord. Augustine refuses to evade the apparent counterexample. A dishonest rich person may grow old comfortably and receive a splendid funeral, while the believer remains poor. If the promise meant a guaranteed redistribution of possessions in this life, that sight would appear to defeat it. Worse, it could make dishonest prosperity look like a pattern the believer should imitate. The question concerns fidelity as well as the interpretation of a sentence.\footnote{Augustine, \work{Enarrationes} 33, second exposition, 13.}

His answer identifies another kind of wealth: justice, truth, charity, faith and endurance are goods of the heart, and Christ himself is living bread. They cannot be bought in the market. A costly bed does not confer them; poverty cannot by itself remove them. Augustine does not ask the hungry person to call hunger pleasant. He asks that possessions cease to decide whether fidelity is worthwhile. Need never turns fraud into goodness, and profitable wrongdoing never proves divine approval.\footnote{Augustine, \work{Enarrationes} 33, second exposition, 14.}

Thomas connects this psalm to Paul's trained contentment. He distinguishes spiritual good from perishable possessions and admits deprivation within providence. His physician analogy concerns God's ordering of the whole person toward salvation. It gives no observer grounds to diagnose why this particular neighbor is hungry. The neighbor's duty comes into focus in Thomas's Philippians exposition: the virtue of bearing poverty does not make assistance unnecessary. The sufferer's trust and another person's obligation are different questions.\footnote{Thomas, \work{Super Psalmo} 33, nn. 10--11; \work{Super Philippenses} 4, lectio 2, paragraph 87844.}

Chrysostom follows the turn in Paul's own sentences. Just when his account of contentment might make the donors think their gift superfluous, Paul commends their fellowship. The first turn restrains the benefactor's pride; the second sustains his generosity. Christ's strength frees Paul from being governed by circumstances or owned by patrons. It does not remove him from relationship. Help can be gratefully received without making its giver the source of salvation.\footnote{Chrysostom, \work{Philippians} XV, ad 4:11--14.}

This balance also changes the giver. A gift supplies a need without entitling the donor to obedience or admiration. Its goodness includes the charity it expresses, not merely the quantity transferred. In the intervening verses, omitted from this Sunday's appointment, Paul names the help sent through Epaphroditus and describes its sacrificial character. Chrysostom insists that money does not purchase heavenly things; the merciful purpose and freedom from attachment matter. A small gift can therefore carry a generosity that a price scale misses.\footnote{Phil 4:15--18, context; Chrysostom, \work{Philippians} XV, ad 4:15--18.}

Neither side may disappear. If material relief alone matters, giving becomes a transaction stripped of communion. If only the donor's virtue matters, the recipient can become an instrument of someone else's spiritual progress. Paul's gratitude joins real relief and real fellowship. The Prayer over the Offerings, whose subject relates offerings and prayer to heavenly glory, lets finite gifts be understood within that larger end. It does not transform them into an admission fee, or permit liturgical generosity to substitute for ordinary care.

Chrysostom and Thomas give distinct emphases to the final promise of supply. Chrysostom thinks of donors with households, children and limited means: Paul desires necessities for people who have given, not their enrichment. Thomas places complete fulfillment in heavenly glory through Christ. The first secures the worth of present bodies and their needs; the second gives desire a final object beyond accumulated goods. They are complementary answers, rather than a choice between material concern and hope of heaven.\footnote{Chrysostom, \work{Philippians} XV, ad 4:19--20; Thomas, \work{Super Philippenses} 4, lectio 2.}

Here the two whole-formulary readings correct different failures. Gregory's garment challenges a guest who treats religious presence as enough. Augustine's psalm and Chrysostom's Paul challenge a believer who measures God's fidelity by possessions, or a neighbor who uses spiritual contentment to excuse neglect. Both require charity, but charity does different work in each argument: it is the life fitting to communion, and it is the practical participation by which one person's want becomes another person's concern.

\section{How does present food lead beyond death?}
The Shepherd's table is prepared amid enemies. Its confidence does not depend on every adversary having disappeared. Augustine reads the Church's journey through regeneration, right paths and discipline toward mature nourishment and lasting dwelling. Thomas finds the tables of divine instruction, sacrament and heaven: teaching strengthens against temptation, the altar sustains the faithful, and the kingdom completes their hope. Provision is present within the valley and ordered beyond it.\footnote{Augustine, \work{Enarrationes} 22.1--6; Thomas, \work{Super Psalmo} 22, nn. 1--2.}

Isaiah supplies a measure no prosperous interval can meet: death itself is defeated and tears are wiped away. The surrounding vision includes judgment and refuge for the needy; the appointed passage ends with the Lord's hand resting on the mountain. Its abundance therefore exceeds either private comfort or a momentary escape from trouble. Thomas's relation of the banquet to Christ's Passion gives the present nourishment its depth. The saving gift comes through one who did not avoid suffering, and the feast's completion promises more than a temporary relief from it.\footnote{Is 24:23; 25:1--12, with 25:6--10a appointed; Thomas, \work{Expositio super Isaiam}, cap. 25.}

Grief can consequently remain truthful. A promise that tears will end does not require mourners to pretend they have ended already. The psalm's table and Paul's sharing of distress place divine presence and real pain together. The community can help now without claiming to accomplish what only the final feast gives. This is the second interpretation's distinctive hope: nourishment sustains fidelity and mutual care during an interval whose wounds are neither final nor unreal.

The acclamation asks for sight adequate to that hope. Chrysostom relates the calling to the divine power manifested in Christ's Resurrection and exaltation; Thomas distinguishes the hope of the journey from the stable inheritance of the life to come. Illumination changes what abundance means. Without it, a promised inheritance might be imagined as an endless extension of present buying and possessing. With it, the desired good is communion in the risen Christ.\footnote{Chrysostom, \work{Homilies on Ephesians} III, ad 1:15--22; Thomas, \work{Super Ephesios} 1, lectio 6, paragraph 87792.}

The Johannine Communion option states future likeness through sight. Augustine begins with the verse's present adoption: the children already belong to God, though their fullness has not appeared. Desire stretches their capacity; purification removes resistance to the gift. The interval is therefore active, not merely empty waiting. This account also deepens Gregory's charity: love is the present life of communion growing toward an unveiled fulfillment, rather than a moral entrance examination detached from the feast itself.\footnote{Augustine, \work{First Epistle of John} IV.5--7.}

Likeness does not mean equality with God. Augustine distinguishes the creature's participation from the Son's equality with the Father. The Prayer after Communion joins nourishment with participation in divine life, and this distinction protects the greatness of that gift without confusing giver and recipient. Thomas's several banquets preserve another distinction: present sacramental communion is real, yet does not erase the difference between pilgrimage and heaven. The feast gives more than an outward badge and less than a warrant to announce that the journey is complete.\footnote{Augustine, \work{First Epistle of John} IV.8--9; Thomas, \work{Expositio super Isaiam}, cap. 25.}

The four senses emerge differently from the two questions. Literally, a king calls and inspects in the long Gospel, while Paul thanks helpers during real hardship. Allegorically, Gregory's mixed Church is joined to the Son in charity, while Augustine and Thomas describe Christ nourishing that body on its journey. Morally, the first question searches for love beneath religious presence; the second calls for fidelity in want and generosity in plenty. Anagogically, the first insists on perseverance toward judgment and vision, the second on the promised end of death beyond every present supply.

With the psalm Communion, the good of seeking the Lord has the last scriptural emphasis; with John, it is the sight toward which desire grows. Neither permits a guest to keep the gift while refusing its life. The Lord who pardons also guides; the Lord who nourishes also draws his people toward a communion whose fullness they cannot manufacture. Their answer is love made concrete now, and hope whose final object is God himself.
